\documentclass[12pt]{article}
\usepackage{amsmath,amssymb}

\begin{document}
\section*{{2018 AIME I Problems/Problem 14}}
Source: AoPS Wiki

\subsection*{Solution}
This is easily solved by recursion/dynamic programming. To simplify the problem somewhat, let us imagine the seven vertices on a line $E \leftrightarrow P_4 \leftrightarrow P_5 \leftrightarrow S \leftrightarrow P_1 \leftrightarrow P_2 \leftrightarrow P_3 \leftrightarrow E$ . We can count the number of left/right (L/R) paths of length $\le 11$ that start at $S$ and end at either $P_4$ or $P_3$ . We can visualize the paths using the common grid counting method by starting at the origin $(0,0)$ , so that a right (R) move corresponds to moving 1 in the positive $x$ direction, and a left (L) move corresponds to moving 1 in the positive $y$ direction. Because we don't want to move more than 2 units left or more than 3 units right, our path must not cross the lines $y = x+2$ or $y = x-3$ . Letting $p(x,y)$ be the number of such paths from $(0,0)$ to $(x,y)$ under these constraints, we have the following base cases: $p(x,0) = \begin{cases} 1 & x \le 3 \\ 0 & x > 3 \end{cases} \qquad p(0,y) = \begin{cases} 1 & y \le 2 \\ 0 & y > 2 \end{cases}$ and recursive step $p(x,y) = p(x-1,y) + p(x,y-1)$ for $x,y \ge 1$ . The filled in grid will look something like this, where the lower-left $1$ corresponds to the origin: $\begin{tabular}{|c|c|c|c|c|c|c|c|} \hline 0 & 0 & 0 & 0 & \textbf{89} & & & \\ \hline 0 & 0 & 0 & \textbf{28} & 89 & & & \\ \hline 0 & 0 & \textbf{9} & 28 & 61 & 108 & 155 & \textbf{155} \\ \hline 0 & \textbf{3} & 9 & 19 & 33 & 47 & \textbf{47} & 0 \\ \hline \textbf{1} & 3 & 6 & 10 & 14 & \textbf{14} & 0 & 0 \\ \hline 1 & 2 & 3 & 4 & \textbf{4} & 0 & 0 & 0 \\ \hline 1 & 1 & 1 & \textbf{1} & 0 & 0 & 0 & 0 \\ \hline \end{tabular}$ The bolded numbers on the top diagonal represent the number of paths from $S$ to $P_4$ in 2, 4, 6, 8, 10 moves, and the numbers on the bottom diagonal represent the number of paths from $S$ to $P_3$ in 3, 5, 7, 9, 11 moves. We don't care about the blank entries or entries above the line $x+y = 11$ . The total number of ways is $1+3+9+28+89+1+4+14+47+155 = \boxed{351}$ . (Solution by scrabbler94)

\end{document}
